% Section 4.4, from lectures/L04/appendix/b_exercises.md. Both sets are code; their solution is in
% the repository.

\section{Exercises}
\label{c4:sec:exercises}\label{c4:app:b}

\begin{aside}
\textbf{How to work these.} Both sets extend one program, the \code{ml} codebase you started in
\secref{c3:sec:exercises}, already carried forward into the lecture's exercises directory,
\file{lectures/L04/exercises} in the companion repository, together with a ready Makefile and
\file{main.cpp}. Set~1 declares the neural network class; Set~2 completes it with a full training method.
The network trains with an adaptive learning rate, following the same rule
\code{ml::lin_reg::Adaptive} uses in Chapter~\ref{c2:ch:training}, so no learning rate has to be
chosen anywhere in the program.

\textbf{How to check your work.} The test suite is already in
\file{lectures/L04/exercises/test}, and Exercise~\ref{c4:ex:14} runs it. It's cumulative:
it carries the stub tests from Chapter~\ref{c3:ch:nn} over unchanged and adds component tests for
\code{Shallow}, so it supersedes the suite from \secref{c3:sec:exercises} entirely. It tests the
directory above it by default; \code{make ML_DIR=<path>} points it at a tree of your own elsewhere.

The solution is published in the companion repository after the lecture, and
\file{lectures/L04/README.md} links to it. Try each exercise yourself before looking at it.
\end{aside}

% ------------------------------------------------------------------------------------------
\exerciseset{The Neural Network Class}
\label{c4:set:1}
You'll extend the codebase from Chapter~\ref{c3:ch:nn}, already carried over into this lecture's
exercises directory, with an interface and a class for a simple
neural network consisting of a hidden layer and an output layer, built on top of the
\code{dense_layer::Interface}/\code{Stub} pair you already have.

\codeexercise{Directory structure}
\label{c4:ex:1}
This lecture's exercises directory, \file{lectures/L04/exercises}, is already set up, so there's
nothing to create or copy:

\begin{console}
exercises/
├── include/
│   └── ml/
│       ├── dense_layer/
│       │   ├── interface.hpp    <- from Chapter 3, ready to use
│       │   └── stub.hpp         <- from Chapter 3, ready to use
│       ├── neural_network/
│       │   ├── interface.hpp    <- to implement
│       │   └── shallow.hpp      <- to implement
│       ├── types.hpp            <- from Chapter 3, ready to use
│       └── utils.hpp            <- to implement
├── source/
│   ├── ml/
│   │   ├── neural_network/
│   │   │   └── shallow.cpp      <- to implement
│   │   └── utils.cpp            <- to implement
│   └── main.cpp                 <- ready to use, see Exercise 4.9
├── test/                        <- ready to use, see Exercise 4.14
└── Makefile                     <- ready to use
\end{console}

\noindent The dense layer interface, the stub, and \file{types.hpp} are carried over from
Chapter~\ref{c3:ch:nn} unchanged. Each file you're to implement holds a \code{@todo} comment saying
what goes in it; replace it with your code. Note that the sources now sit under \file{source/ml/},
mirroring \file{include/ml/}. The makefile already compiles
\file{source/ml/neural_network/shallow.cpp} and \file{source/ml/utils.cpp} along with
\file{main.cpp}.

\codeexercise{Utility functions}
\label{c4:ex:2}
Some helpers don't belong to any one class. The first of them is needed in this chapter, and
Chapter~\ref{c5:ch:dense} adds three more to the same pair of files.

\task{The header (\file{ml/utils.hpp})}
In the namespace \code{ml}, declare a single function:
\begin{itemize}
  \item \textbf{\code{initRandGen()}:}
  \begin{itemize}
    \item Initializes the random number generator, once per program run.
    \item Takes no arguments and returns nothing.
    \item Should be marked \code{noexcept}.
  \end{itemize}
\end{itemize}

\task{The implementation (\file{source/ml/utils.cpp})}
Define \code{initRandGen()} in the same namespace:
\begin{itemize}
  \item Add a static local variable named \code{initialized} with an initial value of
    \code{false}. Being static, it keeps its value between calls, which is what makes the seeding
    happen only once.
  \item If \code{initialized} is already \code{true}, return early.
  \item Seed the generator with the current time by calling
    \code{std::srand(std::time(nullptr))}. Include \code{<cstdlib>} and \code{<ctime>} for these.
  \item Set \code{initialized} to \code{true} so the next call does nothing.
\end{itemize}

\textbf{Why a shared file rather than an anonymous namespace?} In Chapter~\ref{c2:ch:training} the
same function started out in an anonymous namespace inside \file{fixed.cpp}, which was fine while
one class needed it, and moved to a shared file the moment \code{Adaptive} needed it too. An
anonymous namespace doesn't scale: it gives every \file{.cpp} its own private copy, each with its
own \code{initialized} flag, so the generator gets reseeded once per file that asks for it.
Seeding must happen once per \emph{program}. A single definition in \file{utils.cpp}, shared
through \file{utils.hpp}, gives exactly that.

\codeexercise{Interface for neural networks}
\label{c4:ex:3}
In the header file \file{ml/neural_network/interface.hpp}, add a namespace named
\code{ml::neural_network}. In this namespace, implement an interface named \code{Interface}:
\begin{itemize}
  \item \textbf{\code{~Interface()}:} should be set to \code{default} and marked \code{virtual} and
    \code{noexcept}.
  \item \textbf{\code{predict(input)}:} pure virtual. \code{input}: read-only floating-point vector
    of the input to base the prediction on. Returns a reference to a floating-point vector holding
    the predicted value. Should be marked \code{[[nodiscard]]} and \code{noexcept} (\textbf{not}
    \code{const}, since the layers' output is updated on every prediction).
\end{itemize}

\codeexercise{The Shallow class: declaration}
\label{c4:ex:4}
In the header file \file{ml/neural_network/shallow.hpp}, add the namespace
\code{ml::neural_network}. Implement a subclass named \code{Shallow} that inherits \code{Interface}
via public inheritance. The class should be marked \code{final}.
\begin{itemize}
  \item \textbf{\code{Shallow()}:} takes \code{hiddenLayer} and \code{outputLayer} (the network's
    hidden layer and output layer, \code{ml::dense_layer::Interface&}), plus \code{trainInput} and
    \code{trainOutput} (read-only, two-dimensional floating-point vectors holding the training
    data's inputs and outputs). Should be marked \code{explicit} and \code{noexcept}.
  \item \textbf{\code{~Shallow()}:} should be marked \code{default}, \code{noexcept}, and
    \code{override}.
  \item \textbf{\code{predict()}:} overrides the corresponding method in the interface. Should be
    marked \code{[[nodiscard]]}, \code{noexcept}, and \code{override}. Attributes aren't
    inherited, so the override needs its own \code{[[nodiscard]]} (see Chapter~\ref{c3:ch:nn}).
  \item \textbf{\code{train(epochCount, precisionThreshold = 0.999999)}:} trains the network
    (implemented in full in Set~2). \code{epochCount}: number of epochs to train (unsigned
    integer). \code{precisionThreshold}: precision at which training stops early (floating-point
    number). Default value: \code{0.999999} (99.9999\%). Returns \code{true} if training was
    carried out, \code{false} otherwise. Should be marked \code{noexcept}.
\end{itemize}
There's no \code{learningRate} argument. The network picks its own rate while training and revises
it as it goes, exactly as \code{ml::lin_reg::Adaptive} does in Chapter~\ref{c2:ch:training}, so
there's nothing for the caller to guess at. The rule itself is written in
Exercise~\ref{c4:ex:11}.

The class should also have the following private methods, used to train in a randomized order:
\begin{itemize}
  \item \textbf{\code{initTrainOrder(setCount)}:} fills \code{myTrainOrder} (see below) with the
    indices \code{0, 1, 2 ... N-1}, where \code{N} is the number of complete training sets.
    \code{setCount}: that set count (unsigned integer). Returns nothing, and should be marked
    \code{noexcept}.
  \item \textbf{\code{randomizeTrainOrder()}:} shuffles the contents of \code{myTrainOrder} into a
    random order. For each index \code{i}, pick a random index \code{r} and swap
    \code{myTrainOrder[i]} and \code{myTrainOrder[r]}. Takes no arguments, returns nothing, and
    should be marked \code{noexcept}.
\end{itemize}
Each training set passes through the network in three steps, feedforward, backpropagation and
optimization (see \secref{c4:sec:training}), and each step drives both layers. Give every step a
private method of its own, so that \code{train()} reads as the three steps rather than the six
layer calls behind them:
\begin{itemize}
  \item \textbf{\code{feedforward(input)}:} feeds \code{input} through the hidden layer and then the
    output layer. \code{input}: read-only floating-point vector. Returns \code{true} if both layers
    accepted their input, \code{false} otherwise. Should be marked \code{noexcept}.
  \item \textbf{\code{backpropagate(reference)}:} computes the error in the output layer and then in
    the hidden layer. \code{reference}: read-only floating-point vector holding the reference values
    for the input last fed forward. Returns \code{true} if both layers succeeded, \code{false}
    otherwise. Should be marked \code{noexcept}.
  \item \textbf{\code{optimize(input, learningRate)}:} updates the parameters of the hidden layer
    and then the output layer. \code{input}: read-only floating-point vector holding the input last
    fed forward. \code{learningRate}: the learning rate to use (floating-point number). Returns
    \code{true} if both layers succeeded, \code{false} otherwise. Should be marked \code{noexcept}.
\end{itemize}
\textbf{Why not let \code{train()} call \code{predict()}?} A prediction is worth nothing but its
return value, which is why \code{predict()} is \code{[[nodiscard]]}. \code{train()} doesn't want
that value, only the feedforward pass behind it, so calling \code{predict()} there would discard it
and fail to compile under \code{-Werror}. \code{predict()} also can't say whether that pass
succeeded, since it returns the output rather than a \code{bool}. \code{feedforward()} does both
jobs: \code{train()} calls it and checks the result, and \code{predict()} calls it and returns the
output layer's output. That's also why \code{feedforward()} isn't marked \code{[[nodiscard]]}:
\code{predict()} has no way to pass the result on, so it ignores it.

One more private method is needed by the training method, implemented in Set~2:
\begin{itemize}
  \item \textbf{\code{precision()}:} computes the network's precision over the whole training data.
    Takes no arguments and returns a floating-point number. Should be marked \code{[[nodiscard]]}
    and \code{noexcept}, but \textbf{not} \code{const}, unlike its namesake in
    Chapter~\ref{c2:ch:training}: it calls \code{predict()}, which feeds both layers forward and
    therefore changes their output.
\end{itemize}
The learning rate rule that \code{train()} also needs doesn't belong in the class at all. It uses no
member of \code{Shallow}, so it goes in an anonymous namespace in \file{shallow.cpp}, out of the
header (see Exercise~\ref{c4:ex:11}).

Feel free to add more (private) methods as needed.

\codeexercise{Removed constructors and operators}
\label{c4:ex:5}
Delete the class's default constructor, copy and move constructors, and the corresponding
operators.

\codeexercise{Private member variables}
\label{c4:ex:6}
Add the following private member variables to \code{Shallow}:
\begin{itemize}
  \item \textbf{\code{myTrainOrder}:} the indices of the training sets as unsigned integers
    (\code{ml::MatrixU32}), in the order they'll be trained in. Initialized by the constructor and
    reshuffled once per epoch.
  \item \textbf{\code{myHiddenLayer}:} reference to the network's hidden layer, obtained via the
    constructor.
  \item \textbf{\code{myOutputLayer}:} reference to the network's output layer, obtained via the
    constructor.
  \item \textbf{\code{myTrainInput}:} reference to the training data's input, obtained via the
    constructor.
  \item \textbf{\code{myTrainOutput}:} reference to the training data's output, obtained via the
    constructor.
\end{itemize}
No member is needed for the learning rate either, nor for the previous precision. Both are local
variables inside \code{train()}, handed to \code{updateLearningRate()} by reference, so every call
starts from the same initial rate rather than inheriting whatever the previous call ended on.

No separate member is needed for the training set count. \code{myTrainOrder} holds one index per
training set, so \code{myTrainOrder.size()} is that count, exactly as \code{myTrainOrder} replaced
\code{mySetCount} in Chapter~\ref{c2:ch:training}.

\codeexercise{Constructor and prediction}
\label{c4:ex:7}
Implement the following in \file{source/ml/neural_network/shallow.cpp}:

\task{The constructor}
\begin{itemize}
  \item Initialize all member variables as described above.
  \item Print an error message and call \code{std::terminate()} (from \code{<exception>}) if the
    output layer's weight count doesn't match the hidden layer's node count. The two layers can
    never be connected in that case, and the constructor has no way to return a failure code to the
    caller.
  \item Compute the number of complete training sets, i.e.\ the smaller of
    \code{trainInput.size()} and \code{trainOutput.size()}. Print an error message and call
    \code{std::terminate()} if it's 0, as in Chapters~\ref{c1:ch:linreg}
    and~\ref{c2:ch:training}: a network without training data can never be trained.
  \item Fill \code{myTrainOrder} by calling \code{initTrainOrder()} with that set count.
  \item Call \code{initRandGen()} (from \file{ml/utils.hpp}, see Exercise~\ref{c4:ex:2}), so the
    generator is seeded before the first shuffle. Doing it here rather than inside
    \code{randomizeTrainOrder()} keeps the one-time setup out of a method called once per epoch.
\end{itemize}

\task{The method \code{feedforward()}}
\begin{itemize}
  \item Perform feedforward through the entire network:
  \begin{enumerate}
    \item Call \code{myHiddenLayer.feedforward(input)} with the given input.
    \item Call \code{myOutputLayer.feedforward(myHiddenLayer.output())} with the hidden layer's
      output as input.
  \end{enumerate}
  \item Return \code{false} as soon as either call fails, without feeding the output layer if the
    hidden layer rejected the input. Return \code{true} otherwise.
\end{itemize}

\task{The method \code{predict()}}
\begin{itemize}
  \item Call \code{feedforward(input)}.
  \item Return \code{myOutputLayer.output()} (a reference; no separate storage variable is needed
    in \code{Shallow}).
\end{itemize}

\codeexercise{Training method (placeholder)}
\label{c4:ex:8}
Implement a temporary version of \code{train()} in \file{source/ml/neural_network/shallow.cpp} that
simply returns \code{true}. The full implementation (feedforward, backpropagation, and optimization
for each training set and epoch, plus the precision check that sets the learning rate) follows in
Set~2.

Give the private methods that only training uses temporary bodies too, so the header and the
source file stay in step until Set~2 fills them in:
\begin{itemize}
  \item \code{backpropagate()} and \code{optimize()} return \code{false}, so a method you forget to
    finish reports failure rather than quietly claiming success.
  \item \code{precision()} returns \code{0.0}.
\end{itemize}

\codeexercise{Compiling and testing}
\label{c4:ex:9}
\file{main.cpp} is provided in full. Read it before you run it:
\begin{itemize}
  \item It creates one \code{ml::dense_layer::Stub} instance for the hidden layer (3 nodes with 2
    weights per node) and one for the output layer (1 node with 3 weights per node). The output
    layer's weight count matches the hidden layer's node count, which is what connects the two.
  \item It creates an \code{ml::neural_network::Shallow} instance from these two layers and the
    2-bit XOR pattern as training data.
  \item The helper function \code{trainAndTest()} prints a prediction for each training input,
    trains the network for 100 epochs with the default precision threshold, and prints the
    predictions again. \code{main()} returns \code{0} if that succeeds, and \code{-1} otherwise.
\end{itemize}
A few more helper functions in the anonymous namespace keep \code{main()} short:
\begin{itemize}
  \item \textbf{\code{printSeparator()}:} prints a line of 80 dashes.
  \item \textbf{\code{printMatrix(matrix)}:} prints the values of a vector separated by spaces,
    without a trailing newline.
  \item \textbf{\code{printPredictions(header, network, trainInput)}:} prints \code{header} followed
    by a colon, then one line per training input holding the input and the network's prediction.
    It takes the network as an \code{ml::neural_network::Interface&} rather than a
    \code{Shallow&}: a prediction is all it needs, so it works for any network.
\end{itemize}
\code{trainAndTest()}, on the other hand, takes a \code{Shallow&}, since \code{train()} isn't part
of \code{ml::neural_network::Interface}. Note that \code{train()} is called without a learning rate;
the network picks its own.

Build and run the program with \code{make}. The placeholder \code{train()} from
Exercise~\ref{c4:ex:8} returns \code{true} straight away, so the program runs to the end already.
You should get the following output (the dense layers are still stubs, so the prediction is always
0.5):

\begin{console}
--------------------------------------------------------------------------------
Predictions before training:
Input: 0 0, predicted output: 0.5
Input: 0 1, predicted output: 0.5
Input: 1 0, predicted output: 0.5
Input: 1 1, predicted output: 0.5
--------------------------------------------------------------------------------
Predictions after training:
Input: 0 0, predicted output: 0.5
Input: 0 1, predicted output: 0.5
Input: 1 0, predicted output: 0.5
Input: 1 1, predicted output: 0.5
--------------------------------------------------------------------------------
\end{console}

% ------------------------------------------------------------------------------------------
\exerciseset{Training the Neural Network}
\label{c4:set:2}
You'll complete the class \code{ml::neural_network::Shallow} you just built by implementing a full
training method: first the precision it measures itself by, then the rule that turns that
measurement into a learning rate, and finally the training loop that uses both.

\codeexercise{Precision calculation}
\label{c4:ex:10}
Implement the private method \code{precision()} in \file{source/ml/neural_network/shallow.cpp}. It
measures how close the network currently is to its training data, the same mean absolute error
(MAE) the models in Chapter~\ref{c2:ch:training} measure, computed over vectors rather than single
values.

Start with a helper function named \code{averageAbs()}, in an anonymous namespace at the top of
\file{shallow.cpp}:

\begin{cppcode}
[[nodiscard]] double averageAbs(const Matrix1d& x, const Matrix1d& y) noexcept
\end{cppcode}

\begin{itemize}
  \item Compute the number of comparisons as the smaller of \code{x.size()} and \code{y.size()}.
  \item Return \code{0.0} if there's nothing to compare.
  \item Otherwise, sum \code{std::abs(x[i] - y[i])} for every index \code{i} below that count, and
    return the sum divided by the count, i.e.\ the average absolute difference between the two
    vectors.
\end{itemize}
Include \code{<cmath>} in \file{shallow.cpp} for \code{std::abs}.

Then implement \code{precision()} itself:
\begin{itemize}
  \item Declare an error sum starting at zero.
  \item For each training set:
  \begin{itemize}
    \item Call \code{predict()} with the set's input and keep the returned reference.
    \item Add \code{averageAbs(reference, prediction)} to the error sum, where \code{reference} is
      the set's training output.
  \end{itemize}
  \item Compute the MAE by dividing the error sum by the number of training sets,
    \code{myTrainOrder.size()}.
  \item Return \code{1.0 - MAE}.
\end{itemize}

\textbf{Why average each set first?} Chapter~\ref{c2:ch:training} divides by the number of training
sets, because a linear regression model produces exactly one value per set. Here a set holds one
value per output node, so summing the raw errors and dividing by the set count would report an
error that grows with the width of the output layer. Averaging within each set first keeps every
set's contribution on the same scale, whatever the width, so \code{1.0 - MAE} reads the same for a
network with one output node as for one with ten. Since every set holds the same number of values
(see below), the mean of the per-set means is exactly the average error \emph{per value}.

\textbf{Why every set has the same width.} \code{train()} only calls \code{precision()} at the end
of a complete epoch, and every epoch has already passed each training output through
\code{myOutputLayer.backpropagate()}, which returns \code{false} when the row's length doesn't match
the layer's node count. A row of the wrong length therefore ends training before
\code{precision()} is ever reached, so every reference row and every prediction holds exactly one
value per output node. \code{averageAbs()} still compares only up to the shorter of its two vectors,
so it can never read out of range even when called from somewhere without that guarantee.

\codeexercise{The learning rate rule}
\label{c4:ex:11}
The learning rate is the one training parameter with no sensible default, so the network sets its
own, using the precision it just computed. Add a function named \code{updateLearningRate()} to the
anonymous namespace from Exercise~\ref{c4:ex:10}, as Chapter~\ref{c2:ch:training} does in
\file{adaptive.cpp}:

\begin{cppcode}
void updateLearningRate(double& learningRate, double& prevPrecision,
                        const double currentPrecision) noexcept
\end{cppcode}

\noindent\textbf{Why not a private method?} The rule touches nothing but its arguments, so it has no
business in the class, and declaring it in \file{shallow.hpp} would expose an implementation detail
to every file that includes the header. The anonymous namespace makes it visible in
\file{shallow.cpp} only.

It takes the current learning rate and the previous precision by reference, since it
updates both, and the precision just measured by value. Give each of the five values it needs a
named constant: a maximum learning rate of \code{0.25}, a minimum of \code{0.01}, a smallest
expected improvement of \code{0.1}, a step of \code{0.05}, and a decrease factor of \code{0.5}.

The rule works on the difference between the two precisions, i.e.\ how much the network improved
since the last evaluation:
\begin{itemize}
  \item \textbf{Improved by at least the expected minimum:} leave the learning rate alone. It's
    working.
  \item \textbf{Improved, but by less than that:} raise the learning rate by the step. Progress has
    stalled, so take bigger steps.
  \item \textbf{Didn't improve at all:} halve the learning rate. The network is overshooting, so
    take smaller steps.
\end{itemize}
As in Chapter~\ref{c2:ch:training}, keep the limits out of the branches and apply them once
afterwards with \code{learningRate = std::clamp(learningRate, minLearningRate, maxLearningRate)}.
It needs \code{<algorithm>}, which \file{shallow.cpp} already includes for the constructor's
\code{std::min()}.

Store the precision just measured in \code{prevPrecision} before returning, so the next call has
something to compare against.

\textbf{Why halve rather than subtract?} Chapter~\ref{c2:ch:training} lowers the rate by the same
step it raises it by. Here the rate goes up by a fixed step and down by a fixed factor, a scheme
known as \emph{additive increase, multiplicative decrease} (AIMD), the same rule TCP uses to adjust
how fast it sends data. The cut scales with the rate: from \code{0.25}, repeated halving eases down
through \code{0.125}, \code{0.0625} and so on, whereas subtracting \code{0.05} would take any rate
below \code{0.06} straight to the floor after a single bad evaluation. And since a raise always adds
\code{0.05}, a rate that was halved after one unlucky evaluation has more than recovered after the
next one that shows progress.

\textbf{The minimum differs from Chapter~\ref{c2:ch:training} too}, which uses \code{1e-6}. The
purpose the minimum is given there is that ``a model that keeps missing can still creep towards the
line rather than stopping dead'', and in a network \code{1e-6} doesn't deliver it: the parameter
updates it allows are far too small to measure, so the precision drifts \emph{down} by a fraction
of a millionth from one evaluation to the next, the rule reads that as another failure to improve,
and the rate stays at the floor for the rest of the run. A network that sits on a plateau for
enough evaluations can halve its way down there, and it never learns anything afterwards. With
\code{0.01} the rate recovers on its own, since a network that is still moving improves sooner or
later, and the rule raises the rate again the moment it does.

The initial rate sits at that floor as well, \code{0.01} rather than the \code{0.1} of
Chapter~\ref{c2:ch:training}, so the network starts cautiously and only speeds up once it has shown
progress.

\codeexercise{Training method}
\label{c4:ex:12}
Replace the temporary version of \code{train()} in \file{source/ml/neural_network/shallow.cpp}
with a full implementation.

\task{Input validation}
\begin{itemize}
  \item Return \code{false} if \code{epochCount == 0}, or if \code{precisionThreshold} falls outside
    the range \code{(0.0, 1.0)}:
  \begin{itemize}
    \item A threshold of \code{1.0} or more can never be reached, since the precision is
      \code{1.0 - MAE} and the mean absolute error can't be negative. Training would always run
      the full epoch count.
    \item A threshold of \code{0.0} or less accepts a network whose mean absolute error is
      \code{1.0} or worse, which defeats the purpose of the check.
  \end{itemize}
  \item There's no learning rate to validate, since the caller no longer supplies one. The layers
    still check the rate they're handed on every call to \code{optimize()}, and the rule above
    keeps it between \code{0.01} and \code{0.25}, well inside the \code{(0.0, 1.0)} the layers
    demand.
  \item The training set count doesn't need checking here. The constructor already guaranteed at
    least one complete set, so \code{myTrainOrder} is never empty.
  \item As in Chapters~\ref{c1:ch:linreg} and~\ref{c2:ch:training}, \code{train()} reports invalid
    arguments through its return value. Only a constructor calls \code{std::terminate()}, because it
    has no way to return a failure code to the caller.
\end{itemize}

\task{Starting from scratch}
\begin{itemize}
  \item Once the arguments are known to be valid, and before the first epoch, reset both layers by
    calling \code{myHiddenLayer.initParams()} and \code{myOutputLayer.initParams()}. The method is
    part of \code{ml::dense_layer::Interface} from Chapter~\ref{c3:ch:nn}, so the network can call
    it without knowing which kind of layer it holds: \code{Stub} implements it as an empty method,
    and \code{Dense} implements it in Chapter~\ref{c5:ch:dense} by drawing new values.
  \item Training therefore always starts from freshly initialized parameters. Calling
    \code{train()} a second time trains a new network rather than continuing the old one, which is
    what makes retraining a network that ended up badly trained worth doing at all: with the
    parameters carried over, a second call would resume from the same bad state it just failed in.
  \item This is a deliberate departure from \code{ml::lin_reg::Fixed} in Chapter~\ref{c1:ch:linreg},
    whose \code{train()} continues from the current parameters so that a model can be trained in
    stages. That only makes sense while the caller sets the learning rate. Here the schedule
    restarts from the initial rate on every call, so the rate the network picks would no longer
    match the parameters it's applied to.
  \item Do the reset \textbf{after} the argument checks, so a rejected call leaves a trained network
    exactly as it was.
\end{itemize}

\task{Training}
\begin{itemize}
  \item Give the evaluation interval (\code{100}) and the initial learning rate (\code{0.01}) named
    constants.
  \item Declare the learning rate and the previous precision as local variables, initialized to the
    initial learning rate and \code{0.0}. Being local means every call starts from the initial rate
    again, rather than inheriting whatever the previous call ended on.
  \item Iterate the desired number of epochs with a for loop:
    \code{for (std::size_t epoch{}; epoch < epochCount; ++epoch)}.
  \item At the start of every epoch, call \code{randomizeTrainOrder()}. As in
    Chapter~\ref{c2:ch:training}, reshuffling each epoch keeps the network from learning anything
    from the order the training data happens to be stored in.
  \item For each epoch, iterate through the training sets in the order \code{myTrainOrder} gives
    rather than sequentially, e.g.\ with a range-based for loop over \code{myTrainOrder}. Note that
    the loop variable is the index into the training data, not the counter itself.
  \item For each training set index \code{x}, perform the following three steps:
  \begin{enumerate}
    \item \textbf{Feedforward:} call \code{feedforward(myTrainInput[x])}. This performs feedforward
      through both the hidden layer and the output layer.
    \item \textbf{Backpropagation:} call \code{backpropagate(myTrainOutput[x])}. This computes the
      error in both layers.
    \item \textbf{Optimization:} call \code{optimize(myTrainInput[x], learningRate)}, passing the
      current learning rate. This updates the parameters in both layers.
  \end{enumerate}
\end{itemize}
Each of the three calls returns \code{bool}. Return \code{false} as soon as any of them fails: a
dimension mismatch means the network is wired wrong, and training on from there would only produce
meaningless numbers.

Replace the temporary bodies of \code{backpropagate()} and \code{optimize()} with full
implementations:

\task{The method \code{backpropagate()}}
\begin{enumerate}
  \item Compute the error in the output layer: \code{myOutputLayer.backpropagate(reference)}.
  \item Compute the error in the hidden layer from the output layer's error and weights:\newline
    \code{myHiddenLayer.backpropagate(myOutputLayer)}.
\end{enumerate}
The order is crucial: the hidden layer's error is computed from the output layer's, so the output
layer has to go first.

\task{The method \code{optimize()}}
\begin{enumerate}
  \item Optimize the hidden layer: \code{myHiddenLayer.optimize(input, learningRate)}.
  \item Optimize the output layer based on the hidden layer's output:\newline
    \code{myOutputLayer.optimize(myHiddenLayer.output(), learningRate)}.
\end{enumerate}
As in \code{feedforward()}, each method returns \code{false} as soon as either layer call fails,
without making the second call, and \code{true} otherwise. Each layer call returns \code{bool} (see
Chapter~\ref{c3:ch:nn}).

\task{Evaluating the training progress}
\begin{itemize}
  \item Every hundredth epoch, except the first, compute the precision once and keep it in a local
    variable:
  \begin{itemize}
    \item Stop training and return \code{true} if it reaches or exceeds \code{precisionThreshold},
      printing the achieved precision and the number of epochs it took. Report the number of epochs
      completed, not the loop index, so the first epoch reads as 1 rather than 0.
    \item Otherwise, pass it to \code{updateLearningRate()}.
  \end{itemize}
  \item Chapter~\ref{c2:ch:training} evaluates every tenth epoch, since a linear regression model
    reaches its threshold in tens of epochs. A network needs thousands, so a tenth-epoch evaluation
    would scan the whole training data hundreds of times for a picture that has barely changed.
    Every hundredth keeps that overhead down without letting the learning rate go unexamined for
    long.
  \item Note that the precision is computed once per evaluation and then used twice. Calling
    \code{precision()} a second time inside \code{updateLearningRate()} would feed every training
    set through the network again for a value that's already in hand.
\end{itemize}

\task{Return value}
\code{false} on invalid arguments or a failed layer call, otherwise \code{true} once the threshold
is reached or the epochs run out.

\codeexercise{Compiling and testing}
\label{c4:ex:13}
\file{main.cpp} needs no changes: it already calls \code{train()} between the two rounds of
predictions, and keeps using \code{ml::dense_layer::Stub} for both layers until a concrete dense
layer implementation exists (see Chapter~\ref{c5:ch:dense}).

Build and run the program again with \code{make}. The output is the same as in
Exercise~\ref{c4:ex:9}, only now the full training loop runs behind it.

Don't worry if the network doesn't predict correctly; the dense
layers are still stubs and therefore don't actually train. What matters is that the code compiles
and runs without errors; a concrete implementation of \code{Dense} is added in
Chapter~\ref{c5:ch:dense}.

The early-stop line stays absent for the same reason: the stub reports the same output whatever
it's fed, so the precision sits at whatever that fixed output happens to score and never moves.
Every run trains for the full epoch count. Your \code{train()} isn't wrong; there's simply nothing
there to learn until Chapter~\ref{c5:ch:dense}.

\codeexercise{Running the tests}
\label{c4:ex:14}
An updated test suite is available in \file{lectures/L04/exercises/test}. It's cumulative:
it carries the stub tests from Chapter~\ref{c3:ch:nn} over unchanged and adds component tests for
\code{Shallow} on top, so it supersedes the suite from \secref{c3:sec:exercises} entirely.

Once you've implemented the exercises above, build and run it from the exercises directory:

\begin{shell}
make -C test
\end{shell}

\noindent All 30 test cases should pass. The ones worth reading check that a prediction reads the
output layer \emph{live} rather than from a stored copy, that training performs exactly one
feedforward per training set per epoch, and that the precision is evaluated on every hundredth
epoch and no other. Nothing else can tell a correct training loop from one that runs a single pass:
both line up dimensionally and both return \code{true}.

The test suite's README, \file{lectures/L04/exercises/test/README.md}, has more
information, including the mistakes this suite deliberately can't catch and what it would take to
catch them.
